\section{Primary Deliverable}\label{sec:primary-deliverable}

\subsection{BEACON System Overview}\label{subsec:beacon-system-overview}

BEACON's five-page flow from Landing, to Setup, Simulation, NASA-TLX, and Debrief, reflects the typical lifecycle of an MSAR operation.

The Landing page provides an introduction to the system, and introduces scenario presets, including a tutorial.
The user can then choose to proceed to the Setup page, or to a dedicated Management page.

The Setup page allows users to choose the number of drones, the mission boundary, and the scenario seed.
Here, they can also list the amount of anomalies on the map, such as people in the water, lifeboats, or debris.
They can also tweak sensor settings here.
This page is intended to act as the planning stage and supports experimentation with different mission parameters.

After setup, the user can then navigate to the Simulation page, which is the core of the system.
This page runs the real-time simulation of the drone swarm across the map, presenting the canvas, alerts, metrics, and mission controls.
The user can pull up a control panel from the bottom of their screen to access controls like drone commands, search pattern generation, and more.
Here, they can also enable administrative tools like enabling or disabling Fog-of-War, which hides anomalies until they have been detected by a drone, or enabling comms degradation simulation.

The user can then end the mission to navigate to a page that prompts the user to fill in an optional NASA-TLX questionnaire, to capture their perceived workload during the mission.
Choosing to fill in the questionnaire then leads to the dedicated NASA-TLX page, after which the user is navigated to a debrief page which presents the user with their weighted NASA-TLX score and a summary of performance metrics.
The user can export this data as JSON, or CSV, to allow for further analysis.

In addition to the above, there is a dedicated Management page, which allows users to manage their scenarios, and view past performance data.
Missions completed on the same session and browser are stored in the browser's local storage, and can be viewed here.
Alternatively, users can also import and export their performance data as JSON files.

These JSONs would then be analyzed to determine insights about the system's performance, and possible improvements for the MSAR operations team.

\subsection{Environment and Scenario Modelling}\label{subsec:environment-and-scenario-modelling}

The environment generator creates a bounded rectangular mission space with a configurable width and height.
The generator then populates the environment with a configurable number of anomalies, which can be people in the water, lifeboats, or debris, or false positives.
Each anomaly has its own detection radius, which is the furthest distance in meters at which a drone under ideal conditions could detect it.
Here, environmental effects like the sea state that affect drone speed, sensor performance, and communication range are also modelled.

Pre-defined scenario presets are also included, intended to reflect plausible MSAR scenarios, to allow for easier system testing and evaluation.
These presets include a ``Calm Bay'', ``Rough Sea'', ``Comms Degradation'', and ``High-Stress'' scaling test scenario, each with different parameters to test different aspects of the system.

\subsection{Swarm Behaviour, Coverage Planning, and Communication Degradation modelling}\label{subsec:swarm-behaviour-coverage-planning-and-communication-degradation-modelling}

Swarm behaviour is defined via a weighted Voronoi Tessellation approach, defining cells on the map that are assigned to drones based on their current position and speed in such a manner to optimise search coverage and ensure that drones complete their tasks at roughly the same time.
This is intended to avoid situations where some drones have open capacity while others are overwhelmed, leading to suboptimal search coverage and increased operator workload.

The search area is partitioned into a variety of polygonal ``cells'', each having an area directly proportional to the drone's speed.
This gives each drone an implicit ``territory'' to search.

For each cell, a set of waypoints is generated using a Boustrophedon (Lawnmower) path planning algorithm (\cite{10.1007/978-1-4471-1273-0_32}), to ensure systematic coverage.
This approach provides a predictable and explainable search pattern, despite being less adaptable than more complex reinforcement learning-based approaches like EN-MASCA (\cite{ENMASCA2025})\@.
These paths are then rendered via canvas layers, and the drones are given a set of waypoints to follow sequentially, before returning to the Drone Hub in the center of the map.

The communications model seeks to simulate the effects of communication degradation on the drone swarm, by introducing configurable communication ranges, and a probabilistic model of communication failure based on distance and environmental effects.
Packet loss can occur based on the drone's distance from the Drone Hub, as well as the predefined sea state, and other factors defined in the scenario configuration.
By default, the system uses information based off of~\cite{s21082848}, assuming the drones are operating in a 4G LTE network.

The delay can get larger and larger as the conditions get worse, until a point where the drone is effectively disconnected from the Drone Hub, and must rely on its last known instructions and local decision-making to continue its mission.
Any detections made while disconnected get stored in the drone's local memory as a ``buffer'' of queued messages, which are then sent to the Drone Hub all at once when the drone reconnects.
Packet drops and delays drive alert generation to influence situational awareness and impact workload, similar to real-world MSAR operations where communication issues are common and have significant impacts on mission success.

Together, these models create realistic and controllable perturbations that allow the accurate replication of real-world MSAR challenges.

